% Copyright (c) 2026 by C. Klukas.
% Licensed under the MIT License. See LICENSE for details.

We evaluate four questions: can an instance resume on a different kind
of host; what does migration cost and how much does page reuse save;
how do AOT jobs change computation time; and which replication properties
do the checks establish? Appendix~\ref{app:reproduce} links raw records,
build configurations and scripts. Each experiment retains its original
build information; measurements from different builds are not combined
into one performance claim.

\subsection{Setup and measurement scope}\label{sec:evaluation-setup}

The performance tests use three machines on one local area network (LAN):
a Mac Studio (Mac13,1, arm64, 10 CPUs), a Raspberry Pi~4 and a uConsole
with Compute Module~5 (CM5; both Linux arm64, four CPUs). For the
2026-10-15 movement tests, the Mac used GCC~16.2 RelWithDebInfo with
hardware bounds checks; both Pis used GCC~16.1 Release with software
bounds checks. These checks prevent memory accesses outside the guest's
allowed range. Mobile instances used WAMR~2.4.5's fast interpreter in
sandboxed workers. Jobs compiled in the later tests used explicit
software bounds checks on all three hosts. These configuration differences
mean that timing differences cannot be attributed to CPU hardware alone.

Durations use local monotonic clocks, which do not jump with changes
to system time. We never calculate an interval by subtracting timestamps
from different hosts. The source's measurement ends when it receives the
arrival-hook acknowledgement, including the return network trip.
Individual phase medians do not sum to this total: phases can overlap,
and destination activation includes fetching the module and selecting a
worker as well as restoring state. The move table gives ten samples per
row, with median phase times and the minimum and maximum acknowledgement
times. The ranges describe observed samples; they are not confidence
intervals.

\subsection{Portability and state representation}\label{sec:evaluation-images}

The counter snapshot test restores a 131\,212-byte image (about 128\,KB):
two 64\,KB pages plus 140 bytes for metadata and the mutable global.
The restored instance returns the expected accumulated count while the
original continues independently. This checks a C++ object containing
pointers without custom state-serialization code. It does not test every
guest compiler or WebAssembly extension.

The current performance tests move between macOS and Linux, but both
ends use arm64 CPUs. Evidence for moves between CPU architectures comes
from the earlier M0 milestone, using ABI version~0: one counter module
resumed in four transfers involving macOS arm64, Linux arm64 and macOS
x86-64 under Rosetta~2. This supports the image representation used at
that milestone; it is not a current performance test on native x86.
Later correctness tests added a physical Windows~11 x86-64 host. Those
basic functional tests are not Windows performance measurements.

Separate long-operation regressions on macOS exercised scan progress,
cancellation and authority checks, and a scan with injected filesystem
latency lasting over two minutes while status requests remained responsive.
Memory tests grew a job's linear memory to a declared 2\,GB allowance;
duration tests checked admission of unknown, unlimited and week-long
budgets. They did not run a computation for a week or reserve 2\,GB of
physical RAM. These are functional checks, not additional performance
measurements; their implementations are in
\sourcefile{cpp/tests/test_file_scans.cpp}{the scan tests} and
\sourcefile{cpp/tests/test_jobs.cpp}{the job tests}.

\subsection{Physical migration and page reuse}\label{sec:evaluation-moves}

The 2026-10-15 tests completed 480 moves across all six directed pairs
of the three machines. Each pair used 0, 1, 4 and 16\,MB of added test
data (\emph{ballast}), filled with a uniform nonzero byte value. There
were ten outbound/return pairs per size. Page caches were cleared before
each outbound trial (the \emph{cold} move); module caches remained after
first use. Each move checked state, memory size, move counters and
successful arrival. Table~\ref{tab:eval-live-move} shows Mac-to-Pi results.
The source records retain every direction and sample, phase definitions
and executable identities. In the table, ACK is the elapsed time from
capture until the source receives the arrival-hook acknowledgement.

% results: live_move [layout={compact},test-prefix={mac → pi, },strip-prefix={mac → pi, },columns={test|capture, median|destination activation/restore, median|capture to arrival-hook acknowledgement, median|capture to arrival-hook acknowledgement, minimum|capture to arrival-hook acknowledgement, maximum|pages sent, median},headers={Trial|Capture|Restore|ACK|Min|Max|Pages},widths={0.24|0.11|0.11|0.11|0.1|0.1|0.07},caption={Mac-to-Pi physical moves, ten samples per row. Times in ms: Capture, Restore and ACK are medians; Min and Max are the observed ACK range. Phase columns do not sum to ACK. Pages is the median count sent.},label={tab:eval-live-move}]
\begin{table}[htbp]
\centering\small
\caption{Mac-to-Pi physical moves, ten samples per row. Times in ms: Capture, Restore and ACK are medians; Min and Max are the observed ACK range. Phase columns do not sum to ACK. Pages is the median count sent.}
\label{tab:eval-live-move}
\setlength{\tabcolsep}{3pt}
\begin{tabular}{p{0.24\linewidth}p{0.11\linewidth}p{0.11\linewidth}p{0.11\linewidth}p{0.1\linewidth}p{0.1\linewidth}p{0.07\linewidth}}
\toprule
Trial & Capture & Restore & ACK & Min & Max & Pages \\
\midrule
0 MB ballast, cold pages & 1.98 & 7.45 & 30.75 & 21.55 & 48.1 & 2 \\
0 MB ballast, return & 1.95 & 6.93 & 26.34 & 24.06 & 31.29 & 2 \\
1 MB ballast, cold pages & 17.43 & 33.32 & 85.55 & 83.59 & 149.28 & 18 \\
1 MB ballast, return & 17.83 & 26.71 & 57.93 & 56.68 & 92.95 & 2 \\
4 MB ballast, cold pages & 65.02 & 122.02 & 297.49 & 289.65 & 362.6 & 66 \\
4 MB ballast, return & 65.18 & 100.57 & 193.28 & 186.88 & 246.28 & 2 \\
16 MB ballast, cold pages & 268.96 & 474.23 & 1181.85 & 1122.15 & 1267.97 & 258 \\
16 MB ballast, return & 257.02 & 471.33 & 845.13 & 821.93 & 905.15 & 2 \\
\bottomrule
\end{tabular}
\end{table}
% /results

For 16\,MB of ballast, cold moves sent 258 pages and return moves sent
two, reusing 256 pages. Median acknowledgement time fell from 1181.85 to
845.13\,ms. Capture still took about 260\,ms and destination activation
and restoration about 470\,ms. Reusing pages reduces network traffic,
but the source still scans and hashes memory, and the destination still
creates and restores the instance. Sending fewer bytes therefore need
not reduce move time proportionally. Uniform ballast compresses well;
these results do not measure bulk throughput for incompressible state.

Earlier LAN tests on 2026-10-12 measured mean request times of
1.16--5.16\,ms over 20 pings per directed pair and File Courier transfer
rates of 20.4--39.6\,MB/s for 16--256\,MB artifacts. These measurements
describe the network context. They neither measure migration time nor
compare mobility with remote calls for the same application. We do not
infer that moving code is faster than leaving it in place.

\subsection{Compute jobs}\label{sec:evaluation-compute}

Table~\ref{tab:eval-jobs} compares the same integer job in interpreted
and warm AOT modes on 2026-10-15, using three million iterations per
sample. Returned checksums were compared. Warm execution excludes
compilation and queueing, which were recorded separately. Dividing the
interpreter median by the AOT median gives ratios of approximately 2.3
on the Mac, 6.8 on the Pi and 6.7 on the CM5. These ratios apply to this
workload and configuration, not to arbitrary C++ code or the Pi algorithm.
There is no native-C++ baseline.

% results: live_three_host_jobs [layout={compact},columns={test|mode|samples|median elapsed|elapsed range},headers={Host|Mode|Samples|Median (ms)|Range (ms)},widths={0.12|0.3|0.1|0.13|0.2},caption={Repeated identical integer jobs. Warm AOT excludes compilation; unequal interpreter sample counts reflect repeated collection.},label={tab:eval-jobs}]
\begin{table}[htbp]
\centering\small
\caption{Repeated identical integer jobs. Warm AOT excludes compilation; unequal interpreter sample counts reflect repeated collection.}
\label{tab:eval-jobs}
\setlength{\tabcolsep}{3pt}
\begin{tabular}{p{0.12\linewidth}p{0.3\linewidth}p{0.1\linewidth}p{0.13\linewidth}p{0.2\linewidth}}
\toprule
Host & Mode & Samples & Median (ms) & Range (ms) \\
\midrule
mac & AOT job & 5 & 31.0 & 30.0 to 32.0 \\
mac & interpreter job & 10 & 72.0 & 68.0 to 82.0 \\
pi & AOT job & 5 & 27.0 & 27.0 to 29.0 \\
pi & interpreter job & 75 & 184.0 & 174.0 to 239.0 \\
uconsole & AOT job & 5 & 11.0 & 11.0 to 11.0 \\
uconsole & interpreter job & 40 & 74.0 & 68.0 to 97.0 \\
\bottomrule
\end{tabular}
\end{table}
% /results

The Pi integration test used jobs to compute 100\,000 hexadecimal digits
of $\pi$, with all three hosts contributing. It finished in 69\,s,
including one-second status polling. The decimal-output check verifies
the known first 100 digits, not the entire result. This tests the
integration of mobile coordination, jobs and artifact output, rather
than measuring a controlled AOT speedup. Windows correctness tests used
interpreted jobs only.

\subsection{Replication: model, simulation and physical test}\label{sec:evaluation-zombies}

On 2026-10-16 the checker explored ten two-replica worlds completely and
two three-replica worlds to depth~11, considering at most 11 steps from
the initial state (Table~\ref{tab:zombie-results}).
Each job has one work step and one observation. Each channel holds at
most three messages, or two in the three-replica worlds. Fault budgets
limit crashes, network cuts, message loss and duplication, early
timeouts (for example, suspecting a healthy reachable main), sink
failure, cancellation and redirection.

The model also simplifies executions. A send to an unavailable host or
full channel is lost without consuming the message-loss budget. Identical
queued resends are combined. Once the crash budget is exhausted, running
hosts' current state is treated as durable. These assumptions limit what
the checked properties imply about the deployed system.

% results: zombie_protocol [layout={compact},columns={test|reachable states|steps|longest shortest path|explored to (bounded)},headers={World|States|Transitions|Depth|Bound},test-prefix={zombie protocol, exhaustive: },strip-prefix={zombie protocol, exhaustive: |two replicas: },widths={0.42|0.16|0.16|0.07|0.07},caption={Finite-model coverage. Complete worlds check four safety invariants and two recovery goals; bounded worlds check only the four invariants. Depth counts steps on a shortest path to a state, taking the maximum over states; Bound is the search cutoff.},label={tab:zombie-results}]
\begin{table}[htbp]
\centering\small
\caption{Finite-model coverage. Complete worlds check four safety invariants and two recovery goals; bounded worlds check only the four invariants. Depth counts steps on a shortest path to a state, taking the maximum over states; Bound is the search cutoff.}
\label{tab:zombie-results}
\setlength{\tabcolsep}{3pt}
\begin{tabular}{p{0.42\linewidth}p{0.16\linewidth}p{0.16\linewidth}p{0.07\linewidth}p{0.07\linewidth}}
\toprule
World & States & Transitions & Depth & Bound \\
\midrule
three replicas: crash, cut & 9188336 & 24420885 &  & 11 \\
three replicas: two crashes & 5461027 & 14745614 &  & 11 \\
crash, cancel & 3935759 & 19606519 & 41 &  \\
crash, cut & 12908826 & 50102693 & 47 &  \\
crash, duplicate & 4380410 & 18416904 & 40 &  \\
crash, early timeout & 7807220 & 30116664 & 47 &  \\
crash, loss & 3961354 & 17259715 & 43 &  \\
cut, loss, early timeout & 4000767 & 19983567 & 41 &  \\
one joining later: crash, cut & 13205168 & 51296859 & 49 &  \\
one joining later: crash, early timeout & 8051595 & 31115057 & 49 &  \\
sink down, redirect & 378252 & 1698056 & 30 &  \\
two crashes & 22356397 & 81063681 & 53 &  \\
\bottomrule
\end{tabular}
\end{table}
% /results

Only three of the ten two-replica worlds allow an early timeout, at most
once each. In the others, the relevant failure-detection timers are
accurate. Recovery goals are not checked on searches stopped at a depth
limit. The results therefore do not prove correctness under arbitrary
timing, unbounded channels or faults, or complete three-replica
executions. The full suite also requires every declared role transition
to be reached somewhere in the tests.

A deterministic simulator adds virtual clocks, message delays and
checkpoint intervals to the same protocol function. In 10\,000 recorded
three-replica runs with random fault schedules, it selected 119\,289
crash or network-cut events and observed 3370 leadership takeovers.
Some selected events had no effect. The longest measured takeover
without a further disturbance was 18.9\,s, below the configured 25.2\,s
bound. The result-journal check covers observations saved durably on
hosts that are running when the sink accepts the result. The simulated
sink accepts once by construction; it does not independently test atomic
updates in distributed storage.

The ten-minute Mesh Chronicle test on 2026-10-13 introduced a late join,
a frozen host, restarts, a network partition and sink downtime. It
delivered a merged life log through two takeovers and one partition.
Recorded roles and ranks passed the trace checker. This is one physical
execution with a chosen failure schedule, not a measurement of general
availability or a proof that recovery always completes. The pure protocol
model omits real transport, SDK action scheduling and storage behavior.
The physical test and host regression tests cover some of these gaps.

\subsection{Test coverage and limits of the evidence}\label{sec:evaluation-tests}

The test suite exercises execution inside the host process and in
sandboxed workers. A saved full-suite run passed all 18 Mac and 17
Windows CTest entries. Later focused tests cover changed runtime, mesh,
job and example behavior on Mac, Windows, CM5 and Pi. The latest
lifecycle/checkpoint tests passed 136 matched cases across both execution
modes and all four hosts. Each run has its own version record; we do not
claim that every later revision reran every earlier test. Continuous
integration (CI) configurations for detecting memory and concurrency
errors, and decoder tests with malformed inputs, provide further
engineering checks. They do not prove that races or vulnerabilities
are absent.

The evidence is limited by the small LAN, different host configurations,
different historical builds, selected workloads and bounded fault
models. We did not compare mobile coordination with equivalent remote
calls, the Python edition or native C++. Performance tests used at most
three machines and did not exercise real NAT or relay deployments.
The results demonstrate feasibility and specific costs. Comparative
benefit and behavior at larger scales remain unestablished.
